%!TEX root = ../report.tex
\documentclass[report.tex]{subfiles}
\begin{document}

\chapter{Methodology}
\label{ch:methodology}

\section{Datasets and Label Scheme}
\label{sec:data}

All adaptive-inference experiments use the UAVid v1.5
release~\cite{uavid2020}, whose properties are listed in
Table~\ref{tab:uavid}. Ruralscapes is used only for the flow-backend
comparison of Section~\ref{sec:backends}.

\begin{table}[t]
    \centering
    \caption[UAVid properties as verified on disk]{UAVid v1.5 properties,
    verified from the release on disk by file inspection.}
    \label{tab:uavid}
    \footnotesize
    \begin{tabular}{@{}ll@{}}
        \toprule
        \textbf{Property} & \textbf{Value} \\
        \midrule
        Sequences                & 42 (20 train / 7 validation / 15 test) \\
        Frames per sequence      & 901, exactly \\
        Frame rate               & \SI{20}{FPS}, exactly \\
        Native resolution        & $3840 \times 2160$ (\SI{4}{K}) \\
        Densely labelled frames  & ten per sequence, one every 100 frames \\
        \bottomrule
    \end{tabular}
\end{table}

Training uses the ten labelled keyframes of each of the 20 training sequences,
200 images. The adaptive pipeline decodes each full video sequentially and is
scored wherever a labelled frame exists.

\label{sec:labels}
UAVid's colour-coded label images are decoded to six classes plus an ignore
class (Table~\ref{tab:classes}). Two consequences matter downstream: moving
and static cars are \emph{merged} into one \texttt{vehicle} class, and
background becomes \texttt{ignore}, excluded from the loss and from every mIoU
reported here.

\begin{table}[t]
    \centering
    \caption[The six-class label scheme]{The six-class scheme used throughout.}
    \label{tab:classes}
    \footnotesize
    \begin{tabular}{@{}rll@{}}
        \toprule
        \textbf{Id} & \textbf{Class} & \textbf{Note} \\
        \midrule
        0   & road       & \\
        1   & vegetation & low vegetation \\
        2   & tree       & \\
        3   & person     & smallest, rarest object \\
        4   & vehicle    & moving and static car, merged \\
        5   & building   & \\
        255 & ignore     & excluded from loss and mIoU \\
        \bottomrule
    \end{tabular}
\end{table}

\section{Quality Metric}
\label{sec:metric}

Quality is mean intersection-over-union, from a confusion matrix accumulated
over all non-ignored pixels:
%
\begin{equation}
    \mathrm{IoU}_c = \frac{\mathrm{TP}_c}{\mathrm{TP}_c + \mathrm{FP}_c + \mathrm{FN}_c},
    \qquad
    \mathrm{mIoU} = \frac{1}{|\mathcal{C}_{\text{gt}}|} \sum_{c \in \mathcal{C}_{\text{gt}}} \mathrm{IoU}_c ,
\end{equation}
%
where $\mathcal{C}_{\text{gt}}$ is the set of classes present in that frame's
ground truth.

\paragraph{Equal class weighting.} The mean is unweighted, so a class
occupying a few hundred pixels counts as much as one occupying a million.
\texttt{person} regions in UAVid are small, so a few misassigned boundary
pixels cost a large \emph{fraction} of that class's IoU and move the aggregate
mIoU as much as a far larger error on \texttt{building} would. This explains
most of the degradation pattern of Section~\ref{sec:whereloss}.

\paragraph{Two masking modes.} mIoU is computed over two different sets of
pixels. The \emph{all-pixel} figure scores every pixel that is not
\texttt{ignore}. The \emph{valid-only} figure also leaves out the pixels the
forward-backward check has rejected, so it measures the warp only where flow
is trustworthy. Every hybrid-against-baseline comparison uses the all-pixel
figure, so both methods are judged on exactly the same pixels; valid-only
numbers appear only in the flow-backend comparison.

\paragraph{One caveat across chapters.} Training-time validation mIoU
(Table~\ref{tab:backbones}) is computed on a $1024^2$ \emph{centre crop}, the
training pipeline's validation transform, whereas the pipeline experiments
score full-resolution frames. The two are not the same measurement, and the gap
between ConvNeXt-Large's \MIoUConvnextLarge{} and the pipeline baseline's
\BaseMIoU{} is mostly that difference rather than a degraded baseline.

\section{Timing and Experimental Design}
\label{sec:timing}

Every timing is wall-clock and is taken with the GPU synchronised, so no work
can hide behind an asynchronous kernel launch. Each frame is charged what the
pipeline spends on it: model time on a model frame, and flow, warp and
consistency-check time on a propagated one. Baseline and hybrid run over the
same video on the same device, so the comparison is paired.

All experiments ran on a single NVIDIA A100 \SI{80}{GB} under PyTorch 2.6
with CUDA 12.4, in \texttt{bfloat16} at native resolution, with the SEA-RAFT
\texttt{spring-L} checkpoint at 12 refinement iterations.

Four experiments follow: the backbone comparison
(Section~\ref{sec:backbones}), the flow-backend comparison
(Section~\ref{sec:backends}), a sweep of $\tau$ over
$\{0.75, 0.90, 0.95, 1.00\}$ (Section~\ref{sec:sweep}), and the full
validation split at $\tau = 0.95$ (Section~\ref{sec:valresults}). Quality is
measured at the \NumGT{} dense frames and timing over all \NumFrames{};
Section~\ref{sec:threats} states the limits of this protocol.

\end{document}
